\clearpage\section{Background and related work}\label{sec:background}
\subsection{Microlearning and instructional design}
Microlearning organises learning into focused units used in short episodes\citep{microlearning}. Brevity alone does not make a unit effective: it needs an attainable outcome, an explanation and a chance to use the idea. GENAILMS therefore defines a module by one terminal objective and its aligned parts, not by a duration. The revised Bloom taxonomy separates knowledge dimensions from cognitive processes\citep{bloom}; a generated Bloom label is an intention, and whether a question really asks for application has to be judged from the question itself. Constructive alignment requires the objective, activities and assessment to agree\citep{alignment}, which gives a reason to reject a question that is correct but off the objective. Cognitive load theory motivates leaving out unrelated detail and placing a worked example next to the lesson\citep{cogload}.
\subsection{Question generation and LLM authoring}
Automatic question generation predates large language models. Heilman and Smith generate factual questions by transformation and rank them\citep{heilman}; the review by Kurdi et al. reports weak control of difficulty and feedback and little comparable evaluation\citep{aqgreview}. Language models can draft whole lessons, examples and questions, with factual, pedagogical and ethical risks\citep{llmed}. A fluent lesson can omit prerequisites or state a detail that no source supports. The design therefore separates schema validation, a self-check by the model and an independent evaluation, and keeps the lecturer responsible for what is published\citep{aiedreview}.
\subsection{Grounded generation and citation evaluation}
Retrieval-augmented generation (RAG) gives the model passages retrieved from an external corpus\citep{rag}. A retrieved passage can be on topic without supporting a particular sentence, so relevance and support are measured separately. ALCE evaluates citation quality apart from correctness\citep{alce}, and FActScore splits a long answer into atomic claims and checks each one\citep{factscore}. GENAILMS uses both ideas: citations are checked against the passage they point to, and claims without a citation are checked too. Ragas applies the same claim decomposition to question answering as faithfulness, and adds answer relevancy\citep{ragas}.
\subsection{LMS integration and comparable systems}
Moodle Book and Quiz are the delivery and assessment activities\citep{moodlebook,moodlequiz}. Moodle's AI subsystem connects providers to actions inside Moodle\citep{moodleai}; H5P supports authoring of interactive slides\citep{h5p}; Gemini Notebook supports exploring supplied sources\citep{notebook}. LTI~1.3 lets an external tool be opened from a course with a signed launch\citep{lti}. Table~\ref{tab:related} compares their documented primary workflows.
\begin{table}[H]\centering\small\caption{Comparable systems and the workflow addressed by GENAILMS.}\label{tab:related}
\begin{tabularx}{\linewidth}{|p{3.0cm}|Y|Y|Y|}\hline
System & Primary workflow & Instructional structure & LMS relationship\\\hline
Heilman--Smith & Generate and rank factual questions & Questions only & Not specified\\\hline
Gemini Notebook & Explore supplied sources & No module structure & Not specified\\\hline
H5P & Author interactive slides and questions & Composed by the author & Embedded content\\\hline
Moodle AI subsystem & Provider actions inside Moodle & Placement-level help & Native\\\hline
GENAILMS & Generate, review and publish modules from course documents & One LO, seven parts, cited sources & Book and Quiz publication; LTI tool\\\hline
\end{tabularx}\end{table}
\subsection{Choice of retrieval components}\label{sec:measurement-setup}
The main components were chosen for the midterm report on an earlier corpus of four computing subjects: 48 lecture PDFs, 2,145 pages, 1,214 chunks and 240 questions (120 English, 120 Vietnamese), each written by Claude Haiku from one chunk and checked against a quote from it. Measurements ran on CPU with three to four cores per service. They are kept here because they decided the components; results on the five-subject corpus are in Section~\ref{sec:results}.

Dense retrieval uses Qwen3-Embedding-0.6B\citep{qwen}, served locally, truncated from 1,024 to 768 dimensions and renormalised. Table~\ref{tab:embedding-measurements} compares it with other embedding models.
\begin{table}[H]\centering\small\caption{Embedding models on 240 bilingual questions over 1,214 chunks (midterm corpus).}\label{tab:embedding-measurements}
\begin{tabularx}{\linewidth}{|Y|r|r|r|r|}\hline
Model & R@1 & R@5 & R@10 & MRR\\\hline
Qwen3, no query instruction & 0.750 & 0.967 & 0.983 & 0.851\\\hline
Qwen3, with query instruction & 0.796 & 0.971 & 0.988 & 0.873\\\hline
BGE-M3\citep{bgem3} & 0.704 & 0.917 & 0.946 & 0.797\\\hline
OpenAI text-embedding-3-small\citep{openaiembedding} & 0.588 & 0.892 & 0.942 & 0.717\\\hline
Gemini embedding-001\citep{geminiembedding} & 0.796 & 0.946 & 0.988 & 0.869\\\hline
\end{tabularx}\end{table}
Instructed Qwen3 had the highest MRR in both languages (English 0.881, Vietnamese 0.865) and runs on the project's own server, so no course text leaves it for embedding. Gemini was close but is an external API. On three CPU cores Qwen3 took 0.15 s per query and about 3.3 GiB of memory.

Dense search is combined with BM25\citep{bm25} by reciprocal rank fusion\citep{rrf}, with constant 60, and the 30 fused candidates are reranked by a cross-encoder\citep{crossencoder}. Table~\ref{tab:retrieval-measurements} compares two rerankers on the same candidates.
\begin{table}[H]\centering\small\caption{Hybrid retrieval and reranking on 240 questions (midterm corpus); time is CPU reranking of 30 candidates with four cores.}\label{tab:retrieval-measurements}
\begin{tabularx}{\linewidth}{|Y|r|r|r|r|}\hline
Configuration & R@1 & MRR@10 & Gold in top 8 & Rerank mean\\\hline
Dense & 0.750 & 0.850 & 0.983 & --\\\hline
Dense + BM25 (RRF) & 0.733 & 0.816 & 0.963 & --\\\hline
Hybrid + BGE reranker base\citep{bgereranker} & 0.750 & 0.820 & 0.933 & 6.04 s\\\hline
Hybrid + Jina reranker v2 INT8\citep{jina} & 0.858 & 0.911 & 0.988 & 2.35 s\\\hline
\end{tabularx}\end{table}
Jina improved rank-one retrieval over hybrid search (41 against 11 discordant questions, exact McNemar $p<0.001$) and was 2.6 times faster than BGE, so it became the self-hosted reranker. Its weights are licensed for non-commercial use only. Adaptive selection of the context size\citep{adaptive}, with a floor of eight chunks, kept 239 of 240 source chunks against 237 at a fixed eight ($p=0.5$); it is used because it costs nothing, not because it was shown to help.
\subsection{Document readers}
PyMuPDF was compared with Docling\citep{pymupdf,docling} on three code-heavy decks (142 pages). PyMuPDF read a page in 0.007 to 0.064 s; Docling took 1.0 to 4.6 s per page and, in two decks, merged or reordered code lines, so that a C++ block no longer compiled. The design keeps PyMuPDF with layout rules, OCR for pages without text and the repairs for formulas and figures described in Section~\ref{sec:results-extraction}.
\subsection{Language models and cost}
All model calls go through a LiteLLM gateway\citep{litellm}, which holds the provider keys, maps aliases to models and records which model served each call. Table~\ref{tab:prices} gives the prices used in the cost estimates.
\begin{table}[H]\centering\small\caption{Text-token prices, USD per million tokens, November 2026.}\label{tab:prices}
\begin{tabularx}{\linewidth}{|Y|p{2.2cm}|r|r|}\hline
Model & Alias & Input & Output\\\hline
Claude Haiku 5.5 (prompts up to 100K tokens) & \code{claude-fast} & 0.10 & 0.50\\\hline
Claude Sonnet 5.5 & \code{claude-pro} & 2.00 & 10.00\\\hline
GPT-6 Luna & \code{openai-fast} & 0.10 & 0.50\\\hline
GPT-6.1 Sol & \code{openai-pro} & 2.00 & 10.00\\\hline
Gemini 3.5 Flash-Lite (paid tier; a free tier exists) & \code{gemini-fast} & 0.30 & 2.50\\\hline
GPT-OSS-120B on Groq & \code{groq-pro} & 0.15 & 0.60\\\hline
\end{tabularx}\end{table}
Prices are from the providers' rate tables\citep{anthropicprice,lunaprice,solprice,geminiprice,groqprice}. Measured token use per module is reported in Section~\ref{sec:results-latency}.
